\section{Union-closed families}

In 2022 Justin Gilmer proved a uniform positive lower bound on the largest frequency in a finite union-closed family, by comparing two entropies. In this section logarithms are natural. The entropy $H$ of the rest of the chapter is in bits, so the quantity used here is $(\ln 2)$ times that entropy; every comparison below is homogeneous of degree one and the base cancels in the frequency bound. For a random variable $X$ with finitely many values write $H(X)=\sum_xP(X=x)\ln\frac1{P(X=x)}$, and $h(p)=p\ln\frac1p+(1-p)\ln\frac1{1-p}$. By Proposition~\ref{prop:gibbs}, scaled by $\ln 2$, one has $H(X)\le\ln m$ if $X$ takes at most $m$ values, with equality for the uniform law.

For random variables $X,Y$ with finitely many values, let $H(X\mid Y=y)$ be the entropy of the law of $X$ given $Y=y$, and
$H(X\mid Y)=\sum_yP(Y=y)H(X\mid Y=y)$.

\begin{lemma}\label{conjectures:lem:condent}
\begin{enumerate}[label=(\alph*)]
\item (Chain rule.) $H(X,Y)=H(Y)+H(X\mid Y)$, and $H(X_1,\dots,X_n)=\sum_{i=1}^nH(X_i\mid X_1,\dots,X_{i-1})$.
\item (Conditioning reduces entropy.) $H(X\mid Y,W)\le H(X\mid Y)$. If $Y$ is a function of $W$, then $H(X\mid W)\le H(X\mid Y)$.
\end{enumerate}
\end{lemma}
\begin{proof}
(a) Write $P(X=x,Y=y)=P(Y=y)P(X=x\mid Y=y)$, take $\ln\frac1{\,\cdot\,}$, which turns the product into a sum, multiply by $P(X=x,Y=y)$ and add over
$(x,y)$. The second formula follows by induction, applied to $(X_1,\dots,X_{n-1})$ and $X_n$.

(b) Fix $y$, and let $\lambda_w=P(W=w\mid Y=y)$ and $q_w(x)=P(X=x\mid Y=y,W=w)$. Then $P(X=x\mid Y=y)=\sum_w\lambda_wq_w(x)$, and since
$\eta(t)=t\ln\frac1t$ is concave on $[0,1]$ (its second derivative is $-\frac1t$), Jensen's inequality gives
$\eta\bigl(\sum_w\lambda_wq_w(x)\bigr)\ge\sum_w\lambda_w\eta(q_w(x))$. Adding over $x$: $H(X\mid Y=y)\ge\sum_w\lambda_wH(X\mid Y=y,W=w)$. Multiply by
$P(Y=y)$ and add over $y$. If $Y$ is a function of $W$, then conditioning on $(Y,W)$ is conditioning on $W$.
\end{proof}

Gilmer's idea is to take two independent uniform random sets $A,B$ in $\mathcal F$. Their union is again in $\mathcal F$, so it cannot have more entropy
than $A$. But if all elements are rare, a union of two independent sets is more random than one set, because each element gets two chances. The whole
proof is the comparison of these two entropies, coordinate by coordinate.

\begin{theorem}[Gilmer's method]\label{conjectures:thm:gilmer}
Suppose that $c>1$ is such that $\mathbb E\,h(UV)\ge c\,\mathbb E[U]\,\mathbb E\,h(U)$ for every pair of independent, identically distributed random
variables $U,V$ with finitely many values in $[0,1]$. Then in every finite union-closed family $\mathcal F$ with at least two sets, some element has
frequency at least $1-\frac1c$.
\end{theorem}
\begin{proof}
Let the ground set be $\{1,\dots,n\}$ and identify a set with its indicator vector in $\{0,1\}^n$. Let $A,B$ be independent and uniform on $\mathcal F$,
and $C=A\cup B$, with coordinates $C_i=\max(A_i,B_i)$. Write $A_{<i}=(A_1,\dots,A_{i-1})$, and similarly for $B$ and $C$. Suppose that every element has
frequency at most $p$, and let us show that $c(1-p)\le1$.

For a vector $a\in\{0,1\}^{i-1}$ with $P(A_{<i}=a)>0$, let $u_i(a)=P(A_i=0\mid A_{<i}=a)$, and put $U_i=u_i(A_{<i})$ and $V_i=u_i(B_{<i})$. These are
independent with the same law, and $\mathbb E[U_i]=P(A_i=0)\ge1-p$. By the chain rule,
\[
\ln|\mathcal F|=H(A)=\sum_{i=1}^nH(A_i\mid A_{<i})=\sum_{i=1}^n\mathbb E\,h(U_i).
\]
For $C$, the vector $C_{<i}$ is a function of $(A_{<i},B_{<i})$, so by Lemma~\ref{conjectures:lem:condent}(b),
$H(C_i\mid C_{<i})\ge H(C_i\mid A_{<i},B_{<i})$. Given $A_{<i}=a$ and $B_{<i}=b$, the bits $A_i$ and $B_i$ are independent (because $A$ and $B$ are), and
$C_i=0$ if and only if both are $0$, which has probability $u_i(a)u_i(b)$. So $H(C_i\mid A_{<i},B_{<i})=\mathbb E\,h(U_iV_i)$, and by the hypothesis,
\[
H(C)=\sum_{i=1}^nH(C_i\mid C_{<i})\ge\sum_{i=1}^n\mathbb E\,h(U_iV_i)\ge c\sum_{i=1}^n\mathbb E[U_i]\,\mathbb E\,h(U_i)\ge c(1-p)\ln|\mathcal F|.
\]
But $C$ takes values in $\mathcal F$ because $\mathcal F$ is union-closed, so $H(C)\le\ln|\mathcal F|$ by Gibbs' inequality. As $|\mathcal F|\ge2$,
$\ln|\mathcal F|>0$, and $c(1-p)\le1$, that is $p\ge1-\frac1c$. Taking for $p$ the largest frequency proves the theorem.
\end{proof}

Everything now rests on an inequality between numbers. The weak form $h(xy)\ge x\,h(y)$ is just concavity: $h(xy)=h\bigl(x\cdot y+(1-x)\cdot0\bigr)\ge
x\,h(y)+(1-x)h(0)$. It gives the hypothesis with $c=1$ only, which proves nothing. We need to gain a constant factor over it, and here is a version with a
short proof.

\begin{lemma}\label{conjectures:lem:hxy}
For $x,y\in[0,1]$,
\[
h(xy)\ge\frac1{1+\ln2}\bigl(x\,h(y)+y\,h(x)\bigr).
\]
\end{lemma}
\begin{proof}
If $x$ or $y$ is $0$, both sides are $0$. For $0<t\le1$ put $t=e^{-a}$ with $a\ge0$; then
\[
\frac{h(t)}t=\ln\frac1t+\frac{1-t}t\ln\frac1{1-t}=a+K(a),\qquad K(a)=(e^a-1)\ln\frac1{1-e^{-a}}\quad(K(0)=0).
\]
\emph{A series for $K$.} With $\ln\frac1{1-z}=\sum_{m\ge1}\frac{z^m}m$ for $z=e^{-a}<1$,
\[
K(a)=\sum_{m\ge1}\frac{e^{-(m-1)a}-e^{-ma}}m=1-\sum_{m\ge1}\frac{e^{-ma}}{m(m+1)},
\]
by regrouping the terms of each power of $e^{-a}$ (the series converge absolutely), and this also holds at $a=0$, as $\sum\frac1{m(m+1)}=1$. In
particular $K\ge0$, and for $a,b\ge0$ the \emph{defect} is
\[
D(a,b)=K(a)+K(b)-K(a+b)=\sum_{m\ge1}\frac{(1-e^{-ma})(1-e^{-mb})}{m(m+1)}.
\]
\emph{The defect is at most $(a+b)\ln2$.} The function $\ell(z)=\ln(1-e^{-z})$ is concave on $]0,\infty[$, since $\ell'(z)=\frac1{e^z-1}$ decreases. So
$(1-e^{-ma})(1-e^{-mb})\le(1-e^{-ms})^2$ with $s=\frac{a+b}2$, and $D(a,b)\le\sum_m\frac{(1-e^{-ms})^2}{m(m+1)}=:F(s)$. We show $F(s)\le2s\ln2$. Both sides
vanish at $s=0$; differentiating term by term (legitimate for $s>0$, where the differentiated series converges uniformly near $s$), with $t=e^{-s}$,
\[
F'(s)=2\sum_{m\ge1}\frac{t^m-t^{2m}}{m+1}=2\Bigl(\frac1t\ln\frac1{1-t}-\frac1{t^2}\ln\frac1{1-t^2}\Bigr)=\frac2{t^2}\Bigl(\ln(1+t)+(1-t)\ln(1-t)\Bigr),
\]
using $\sum_{m\ge1}\frac{z^m}{m+1}=\frac1z\ln\frac1{1-z}-1$. Let $Q(t)=t^2\ln2-\ln(1+t)-(1-t)\ln(1-t)$ on $[0,1]$. Then $Q(0)=Q(1)=0$,
$Q'(t)=2t\ln2+\frac t{1+t}+\ln(1-t)$, $Q'(0)=0$, $Q''(t)=2\ln2+\frac1{(1+t)^2}-\frac1{1-t}$, and $Q'''(t)=-\frac2{(1+t)^3}-\frac1{(1-t)^2}<0$. So $Q''$
decreases from $Q''(0)=2\ln2>0$ to $-\infty$: $Q'$ first increases from $0$, then decreases to $-\infty$, so it is positive and then negative; $Q$
increases and then decreases, and since $Q(0)=Q(1)=0$, $Q\ge0$ on $[0,1]$. Hence $F'(s)\le2\ln2$, and $F(s)\le2s\ln2$.

\emph{Conclusion.} Let $x=e^{-a}$, $y=e^{-b}$, and $c'=\frac1{1+\ln2}$, so that $1-c'=c'\ln2$. Dividing the claim by $xy$, we must show
$a+b+K(a+b)\ge c'\bigl(a+K(a)+b+K(b)\bigr)$. The difference of the two sides is
\[
(1-c')(a+b)+(1-c')K(a+b)-c'D(a,b)\ge(1-c')(a+b)-c'(a+b)\ln2=0,
\]
because $K\ge0$ and $D(a,b)\le(a+b)\ln2$.
\end{proof}

\begin{corollary}\label{conjectures:cor:ucconst}
In every finite union-closed family with at least two sets, some element has frequency at least $\frac{1-\ln2}2=0.1534\ldots$
\end{corollary}
\begin{proof}
For $U,V$ independent with the same law, Lemma~\ref{conjectures:lem:hxy} and independence give
\[
\begin{aligned}\mathbb E\,h(UV)&\ge\frac1{1+\ln2}\bigl(\mathbb E[U]\,\mathbb E\,h(V)+\mathbb E[V]\,\mathbb E\,h(U)\bigr)\\&=\frac2{1+\ln2}\,\mathbb E[U]\,\mathbb E\,h(U).\end{aligned}
\]
Apply
Theorem~\ref{conjectures:thm:gilmer} with $c=\frac2{1+\ln2}$: $1-\frac1c=1-\frac{1+\ln2}2=\frac{1-\ln2}2$.
\end{proof}
